% Options for packages loaded elsewhere
\PassOptionsToPackage{unicode}{hyperref}
\PassOptionsToPackage{hyphens}{url}
\documentclass[
  12pt,
]{ctexart}
\usepackage{xcolor}
\usepackage{amsmath,amssymb}
\setcounter{secnumdepth}{-\maxdimen} % remove section numbering
\usepackage{iftex}
\ifPDFTeX
  \usepackage[T1]{fontenc}
  \usepackage[utf8]{inputenc}
  \usepackage{textcomp} % provide euro and other symbols
\else % if luatex or xetex
  \usepackage{unicode-math} % this also loads fontspec
  \defaultfontfeatures{Scale=MatchLowercase}
  \defaultfontfeatures[\rmfamily]{Ligatures=TeX,Scale=1}
\fi
\usepackage{lmodern}
\ifPDFTeX\else
  % xetex/luatex font selection
\fi
% Use upquote if available, for straight quotes in verbatim environments
\IfFileExists{upquote.sty}{\usepackage{upquote}}{}
\IfFileExists{microtype.sty}{% use microtype if available
  \usepackage[]{microtype}
  \UseMicrotypeSet[protrusion]{basicmath} % disable protrusion for tt fonts
}{}
\makeatletter
\@ifundefined{KOMAClassName}{% if non-KOMA class
  \IfFileExists{parskip.sty}{%
    \usepackage{parskip}
  }{% else
    \setlength{\parindent}{0pt}
    \setlength{\parskip}{6pt plus 2pt minus 1pt}}
}{% if KOMA class
  \KOMAoptions{parskip=half}}
\makeatother
\usepackage{longtable,booktabs,array}
\usepackage{calc} % for calculating minipage widths
% Correct order of tables after \paragraph or \subparagraph
\usepackage{etoolbox}
\makeatletter
\patchcmd\longtable{\par}{\if@noskipsec\mbox{}\fi\par}{}{}
\makeatother
% Allow footnotes in longtable head/foot
\IfFileExists{footnotehyper.sty}{\usepackage{footnotehyper}}{\usepackage{footnote}}
\makesavenoteenv{longtable}
\setlength{\emergencystretch}{3em} % prevent overfull lines
\providecommand{\tightlist}{%
  \setlength{\itemsep}{0pt}\setlength{\parskip}{0pt}}
\usepackage{bookmark}
\IfFileExists{xurl.sty}{\usepackage{xurl}}{} % add URL line breaks if available
\urlstyle{same}
\hypersetup{
  hidelinks,
  pdfcreator={LaTeX via pandoc}}

\author{SCX}
\date{2026}

\begin{document}

\section{{[}Archived{]} Proposition 1: Analogy to Arrow's
Impossibility Theorem}\label{archived-proposition-1-analogy}

\begin{quote}
\textbf{Removal Date}: 2026-06-27

\textbf{Reason for Removal}: This section draws an analogy between the insufficiency of global expert ranking in SCX and Arrow's Impossibility Theorem. Although the analogy is intuitive and suggestive, it suffers from the following issues:

\begin{enumerate}
\def\labelenumi{\arabic{enumi}.}
\tightlist
\item
  \textbf{Condition Mismatch}: Arrow's Theorem requires four conditions --- universality (unrestricted domain), Pareto optimality, non-dictatorship, and independence of irrelevant alternatives --- to deduce impossibility. SCX's Proposition 1 only requires that no globally optimal expert exists when there is state-level ranking crossing --- it does not require all four of Arrow's conditions to hold simultaneously.
\item
  \textbf{Misleading}: The analogy suggests that SCX's impossibility is on the same level as Arrow's Theorem (i.e., no global ranking can exist regardless of how finely states are partitioned), but in reality SCX's ``impossibility'' is conditional --- when the state partition is sufficiently fine and free of crossing, a global ranking is possible.
\item
  \textbf{Limited Theoretical Value}: This analogy makes no substantive contribution to the analysis or derivation of the SCX framework; it belongs to heuristic rhetoric rather than mathematical argumentation.
\end{enumerate}
\end{quote}

\begin{center}\rule{0.5\linewidth}{0.5pt}\end{center}

\subsection{4 Connection to Arrow's Impossibility
Theorem}\label{connection-to-arrows-impossibility-theorem}

The global expert ranking problem bears a deep analogy to social choice theory.

\begin{longtable}[]{@{}ll@{}}
\toprule\noalign{}
Arrow's Theorem & SCX \\
\midrule\noalign{}
\endhead
\bottomrule\noalign{}
\endlastfoot
Voter \(i\) & State \(s\) \\
Candidate \(a\) & Expert \(f_m\) \\
Social preference ordering & Global risk ranking \(R_m\) \\
Pareto optimality & State-level ranking consistency \\
Independence of irrelevant alternatives & Independence upon introducing new experts \\
\end{longtable}

Arrow's Impossibility Theorem states: there exists no function mapping individual preferences to social preferences that simultaneously satisfies Pareto optimality, non-dictatorship, and independence of irrelevant alternatives.

SCX's state-level rankings correspond to ``individual preferences'' (each state \(s\) has its own expert ranking), and the global ranking corresponds to the ``social preference.'' Proposition 1 shows: \textbf{the global ranking cannot consistently reflect state-level rankings}, just as in Arrow's Theorem the social preference cannot simultaneously satisfy all rationality conditions.

Unlike Arrow's Theorem, however, SCX's ``impossibility'' does not require universality (unrestricted domain) and only points to the existence of a counterexample.

\begin{center}\rule{0.5\linewidth}{0.5pt}\end{center}

\subsection{Preservation Notes}\label{preservation-notes}

This section has been removed from \texttt{01\_global\_ranking\_insufficiency.md}. If one wishes to cite the analogical relationship between Arrow's Theorem and SCX, please cite Arrow (1951) directly and provide your own argumentation, rather than citing this section.

\end{document}
